%% GENERATED by generate_papers.py -- edit the bank_*.py files, not this file.
\documentclass[11pt,a4paper]{article}
\usepackage[a4paper, top=20mm, bottom=20mm, left=18mm, right=18mm, headheight=14pt]{geometry}
\usepackage[T1]{fontenc}
\usepackage{lmodern}
\usepackage{amsmath, amssymb, mathtools}
\usepackage{braket}
\usepackage{array, tabularx, booktabs, multicol}
\usepackage{enumitem}
\usepackage[dvipsnames]{xcolor}
\usepackage[most]{tcolorbox}
\usepackage{needspace}
\usepackage{fancyhdr, lastpage}
\usepackage{pifont}
\usepackage{tikz}
\usepackage[hidelinks]{hyperref}
\newcolumntype{L}{>{\raggedright\arraybackslash}X}
\definecolor{brand}{HTML}{1F4E79}
\definecolor{okgreen}{HTML}{1B7F3B}
\setlength{\parindent}{0pt}
\setlength{\parskip}{3pt}
\newcommand{\optlabel}[1]{\textbf{(#1)}}
\newcommand{\correctlabel}[1]{\textcolor{okgreen}{\textbf{(#1)}\,\ding{51}}}
\newenvironment{qitem}[1]{\begin{list}{}{\setlength{\leftmargin}{3.3em}\setlength{\labelwidth}{3.1em}%
  \setlength{\labelsep}{0.2em}\setlength{\topsep}{3pt}\setlength{\parsep}{3pt}}\item[\textbf{Q.#1}\hfill]}{\end{list}}
\newcommand{\qsep}{\par\vspace{2pt}{\color{black!15}\hrule}\vspace{4pt}}
\newcommand{\qmeta}[2]{{\footnotesize\color{black!60}\textsc{#1} \;$\cdot$\; #2}\par}
\newtcolorbox{passagebox}[1]{enhanced, breakable, colback=brand!4, colframe=brand!70, fonttitle=\bfseries,
  title={#1}, left=6pt, right=6pt, top=4pt, bottom=4pt}
\newtcolorbox{sectionbanner}{enhanced, colback=brand, colframe=brand, coltext=white, fontupper=\large\bfseries,
  left=8pt, right=8pt, top=6pt, bottom=6pt, sharp corners}
\newtcolorbox{infobox}[1]{enhanced, breakable, colback=black!3, colframe=black!45, fonttitle=\bfseries,
  title={#1}, left=6pt, right=6pt, top=4pt, bottom=4pt}
\newtcolorbox{solbox}[1]{enhanced, breakable, colback=okgreen!3, colframe=okgreen!60!black, fonttitle=\bfseries,
  title={Answer: (#1) \quad Detailed solution}, left=6pt, right=6pt, top=3pt, bottom=4pt, before skip=4pt, after skip=6pt}
\newcommand{\QAOAop}{\mathrm{QAOA}}
\newcommand{\bub}[1]{\tikz[baseline=-0.65ex]\node[draw, circle, inner sep=0pt, minimum size=4.6mm, font=\scriptsize]{#1};}
\allowdisplaybreaks

\begin{document}
\pagestyle{fancy}\fancyhf{}
\fancyhead[L]{\small\color{black!70}Quantum Approximate Optimization Algorithm --- Assessment}
\fancyhead[R]{\small\color{black!70}\textbf{Question Paper}}
\fancyfoot[C]{\small Page \thepage\ of \pageref*{LastPage}}
\renewcommand{\headrulewidth}{0.4pt}

\thispagestyle{empty}
\begin{center}
{\color{brand}\rule{\linewidth}{1.2pt}}\\[6pt]
{\Huge\bfseries\color{brand} Quantum Approximate Optimization Algorithm}\\[4pt]
{\Large Assessment of the QAOA algorithm (Max-Cut)}\\[4pt]
{\large\bfseries QUESTION PAPER}\\[2pt]
{\color{brand}\rule{\linewidth}{1.2pt}}\\[1pt]
\end{center}
\vspace{-2pt}
\begin{tabularx}{\linewidth}{@{}lXlr@{}}
\textbf{Time allowed:} & 3 hours & \textbf{Maximum marks:} & 240\\
\textbf{Questions:} & 80 (4 sections $\times$ 20) & \textbf{Answer key:} & separate Answer Book
\end{tabularx}
\vspace{6pt}

\textbf{Name of the candidate:} \hrulefill\quad \textbf{Roll No.:} \rule{35mm}{0.4pt}

\vspace{6pt}
\begin{infobox}{General instructions}
\begin{enumerate}[leftmargin=1.4em, itemsep=1pt, topsep=1pt]
\item This question paper has \textbf{four sections} of \textbf{20 questions} each:
\textbf{A}~Mathematics and Graph Theory for QAOA, \textbf{B}~Quantum Mechanics, \textbf{C}~QAOA Basics, \textbf{D}~Advanced QAOA and Implementation.
\item Every question is of the \textbf{single-correct} type: four options, exactly one correct. Question formats used:
single-correct MCQ, matching lists (List-I/List-II), Assertion--Reason, Statement~I/II and paragraph-based questions.
\item Mark your answers on the \textbf{Response Sheet} at the end of this paper by darkening ONE bubble per question.
\item Calculators are \textbf{not} allowed. Rough work may be done on separate sheets.
\item Unless a question says otherwise, use the \textbf{conventions} printed below; they apply to every section.
\item The detailed solutions are in the separate \emph{Answer Book}; do not consult it before finishing the test.
\end{enumerate}
\end{infobox}
\begin{infobox}{Marking scheme (every question)}
Each question has \textbf{four} options (A), (B), (C) and (D). \textbf{ONLY ONE} of these four options is correct.\\[2pt]
\begin{tabular}{@{}ll@{}}
\emph{Full marks}: & $+3$ if ONLY the correct option is chosen;\\
\emph{Zero marks}: & $\phantom{+}0$ if none of the options is chosen (question unanswered);\\
\emph{Negative marks}: & $-1$ in all other cases (wrong option, or more than one option chosen).
\end{tabular}
\end{infobox}

\begin{infobox}{Conventions and notation used throughout the paper}
\small
\begin{itemize}[leftmargin=1.2em, itemsep=1pt, topsep=2pt]
\item Graph $G = (V, E)$ with nodes $0, 1, \dots, n-1$ and edge weights $w_{ij}$ ($w_{ij} = 1$ unless stated); $W = \sum_{(i,j)\in E} w_{ij}$.
      A split is a bitstring $z = z_0z_1\dots z_{n-1}$ ($z_i \in \{0,1\}$ = group of node $i$, \textbf{node 0 written first});
      $\bar z$ is $z$ with every bit flipped. Spins: $s_i = (-1)^{z_i} = 1 - 2z_i$, the eigenvalue of $Z_i$ on $\ket z$.
\item Cut value $C(z) = \sum_{(i,j)\in E} w_{ij}\,[z_i \neq z_j] = \sum_{(i,j)} w_{ij}\frac{1 - s_is_j}{2}$; $C_{\max} = \max_z C(z)$.
\item \textbf{Notebook convention} (\texttt{QAOA.ipynb}): $H_C = \sum_{(i,j)\in E}\frac{w_{ij}}{2}Z_iZ_j$, so $C = \frac W2 - H_C$;
      $H_M = \sum_i X_i$; $\ket{\boldsymbol\gamma,\boldsymbol\beta} = e^{-i\beta_pH_M}e^{-i\gamma_pH_C}\cdots e^{-i\beta_1H_M}e^{-i\gamma_1H_C}\ket{+}^{\otimes n}$;
      $F_p(\boldsymbol\gamma,\boldsymbol\beta) = \braket{\boldsymbol\gamma,\boldsymbol\beta|C|\boldsymbol\gamma,\boldsymbol\beta}$; approximation ratio $r = F_p/C_{\max}$.
      The angle vector is \texttt{params} $= [\gamma_1, \dots, \gamma_p, \beta_1, \dots, \beta_p]$.
\item Gates: $H = \frac{1}{\sqrt2}\begin{psmallmatrix}1&1\\1&-1\end{psmallmatrix}$; $X, Y, Z$ are the Pauli matrices; $R_X(\theta) = e^{-i\theta X/2}$,
      $R_Z(\theta) = e^{-i\theta Z/2}$; CNOT$(c, t)$ flips qubit $t$ when qubit $c$ is 1. Two-qubit matrices use the basis order
      $\ket{00}, \ket{01}, \ket{10}, \ket{11}$ (lower-numbered qubit first); $\ket{\pm} = \frac{1}{\sqrt2}(\ket0 \pm \ket1)$.
\item \textbf{Notebook defaults:} random graph with 15 nodes and edge probability $0.3$ (\texttt{SEED = 42}: 30 edges), $p = 2$, COBYLA with
      80 evaluations, ramp initial angles, 4000 shots, exact objective via \texttt{cudaq.observe}. A \emph{previously recorded run} means the
      outputs stored in the notebook (CUDA-Q 0.16.0, CPU target \texttt{qpp-cpu}).
\item All measurements are in the computational basis; $\ln$ is the natural logarithm. Numerical answers are rounded.
\end{itemize}
\end{infobox}
\newpage

\begin{sectionbanner}SECTION A \enspace$\cdot$\enspace Mathematics and Graph Theory for QAOA \hfill Q.1--20\end{sectionbanner}

{\small \textbf{Syllabus:} graphs, cuts and Max-Cut, bipartite graphs and odd cycles, complexity, QUBO and Ising forms, binary and spin variables, probability and sampling statistics, matrix exponentials and commutators.\\ \textbf{Formats:} Single-correct MCQ (Q.1--19); Matching lists (Q.2--13); Assertion--Reason (Q.7--20); Statement I / II (Q.8--11); Paragraph based (Q.17--18). Each question: $+3$ / $0$ / $-1$.}
\par\medskip

\needspace{7\baselineskip}
\begin{qitem}{1}
A graph has $n = 15$ nodes. The number of bitstrings $z \in \{0,1\}^{15}$, and the number of \emph{distinct}
partitions $\{S, V\setminus S\}$ of the node set (unordered; $S = \varnothing$ allowed) that these bitstrings describe, are
\par\smallskip\noindent\begin{tabularx}{\linewidth}{@{}*{4}{L}@{}}
\optlabel{A}\,$32\,768$ and $16\,384$ & \optlabel{B}\,$16\,384$ and $8\,192$ & \optlabel{C}\,$32\,768$ and $16\,383$ & \optlabel{D}\,$32\,768$ and $32\,768$
\end{tabularx}
\end{qitem}
\qsep

\needspace{7\baselineskip}
\begin{qitem}{2}
Match each graph (List-I) with its maximum cut $C_{\max}$ (List-II). All edges have weight 1.
\par\smallskip\noindent{\renewcommand{\arraystretch}{1.35}%
\begin{tabularx}{\linewidth}{@{}c>{\raggedright\arraybackslash}p{0.33\linewidth}@{\hspace{1.2em}}cL@{}}
\toprule
\multicolumn{2}{@{}l}{\textbf{List-I}} & \multicolumn{2}{l@{}}{\textbf{List-II}} \\ \midrule
\textbf{(P)} & cycle $C_9$ (9 nodes, 9 edges) & \textbf{(1)} & $12$ \\
\textbf{(Q)} & complete graph $K_5$ & \textbf{(2)} & $9$ \\
\textbf{(R)} & complete bipartite graph $K_{3,3}$ & \textbf{(3)} & $6$ \\
\textbf{(S)} & Petersen graph (10 nodes, 15 edges, 3-regular) & \textbf{(4)} & $8$ \\
\bottomrule
\end{tabularx}}
\par\smallskip\noindent\begin{tabularx}{\linewidth}{@{}*{2}{L}@{}}
\optlabel{A}\,P$\to$4, Q$\to$3, R$\to$2, S$\to$1 & \optlabel{B}\,P$\to$2, Q$\to$3, R$\to$1, S$\to$4 \\[3pt]
\optlabel{C}\,P$\to$4, Q$\to$2, R$\to$3, S$\to$1 & \optlabel{D}\,P$\to$4, Q$\to$3, R$\to$1, S$\to$2
\end{tabularx}
\end{qitem}
\qsep

\needspace{7\baselineskip}
\begin{qitem}{3}
A weighted graph has the edges $(0,1)$ with weight 3, $(1,2)$ with weight 1, $(2,3)$ with weight 2, $(3,0)$ with
weight 4 and $(0,2)$ with weight 2. Every node is put into group 0 or group 1 independently with probability $\tfrac12$.
The expected cut value is
\par\smallskip\noindent\begin{tabularx}{\linewidth}{@{}*{4}{L}@{}}
\optlabel{A}\,$12$ & \optlabel{B}\,$6$ & \optlabel{C}\,$3$ & \optlabel{D}\,$2.4$
\end{tabularx}
\end{qitem}
\qsep

\needspace{7\baselineskip}
\begin{qitem}{4}
For the path graph $0 - 1 - 2$ (two unit edges), the cut value written as a polynomial in the binary variables
$z_0, z_1, z_2 \in \{0, 1\}$ is
\par\smallskip\noindent\begin{tabularx}{\linewidth}{@{}*{2}{L}@{}}
\optlabel{A}\,$z_0 + 2z_1 + z_2 - z_0z_1 - z_1z_2$ & \optlabel{B}\,$z_0 + 2z_1 + z_2 - 2z_0z_1 - 2z_1z_2$ \\[3pt]
\optlabel{C}\,$z_0 + z_1 + z_2 - 2z_0z_1 - 2z_1z_2$ & \optlabel{D}\,$2 - z_0 - 2z_1 - z_2 + 2z_0z_1 + 2z_1z_2$
\end{tabularx}
\end{qitem}
\qsep

\needspace{7\baselineskip}
\begin{qitem}{5}
The notebook's default graph is a random graph on $n = 15$ nodes in which each of the possible edges is present
independently with probability $0.3$. The expected number of edges and its standard deviation are
\par\smallskip\noindent\begin{tabularx}{\linewidth}{@{}*{4}{L}@{}}
\optlabel{A}\,$45$ and $3.97$ & \optlabel{B}\,$31.5$ and $22.05$ & \optlabel{C}\,$63$ and $6.64$ & \optlabel{D}\,$31.5$ and $4.70$
\end{tabularx}
\end{qitem}
\qsep

\needspace{7\baselineskip}
\begin{qitem}{6}
The 15-node graph of the notebook has the degree sequence (nodes $0$ to $14$)
\[ (d_0, \dots, d_{14}) = (2, 3, 5, 4, 5, 4, 6, 1, 4, 3, 4, 2, 4, 6, 7). \]
The number of edges, and the number of nodes of odd degree, are
\par\smallskip\noindent\begin{tabularx}{\linewidth}{@{}*{4}{L}@{}}
\optlabel{A}\,$30$ and $5$ & \optlabel{B}\,$60$ and $6$ & \optlabel{C}\,$30$ and $6$ & \optlabel{D}\,$30$ and $9$
\end{tabularx}
\end{qitem}
\qsep

\needspace{7\baselineskip}
\begin{qitem}{7}
\textbf{Assertion (A):} The cycle $C_7$ (7 nodes, 7 unit edges) has $C_{\max} = 6$.\\[1mm]
\textbf{Reason (R):} Every graph with an odd number of nodes has $C_{\max} < |E|$.
\begin{itemize}[leftmargin=2.6em, labelsep=0.4em, itemsep=2pt, topsep=3pt]
\item[\optlabel{A}] Both Assertion (A) and Reason (R) are true, and R is the correct explanation of A
\item[\optlabel{B}] Both A and R are true, but R is \emph{not} the correct explanation of A
\item[\optlabel{C}] A is true, but R is false
\item[\optlabel{D}] A is false, but R is true
\end{itemize}
\end{qitem}
\qsep

\needspace{7\baselineskip}
\begin{qitem}{8}
\textbf{Statement I:} A polynomial-time algorithm that solves Max-Cut \emph{exactly} on every graph is known.\\[1mm]
\textbf{Statement II:} Finding a \emph{minimum} cut of a graph (both groups non-empty, non-negative weights) is NP-hard.
\begin{itemize}[leftmargin=2.6em, labelsep=0.4em, itemsep=2pt, topsep=3pt]
\item[\optlabel{A}] Both Statement I and Statement II are true
\item[\optlabel{B}] Both Statement I and Statement II are false
\item[\optlabel{C}] Statement I is true, but Statement II is false
\item[\optlabel{D}] Statement I is false, but Statement II is true
\end{itemize}
\end{qitem}
\qsep

\needspace{7\baselineskip}
\begin{qitem}{9}
For the triangle $K_3$ (unit weights) consider $H_C = \tfrac12\big(Z_0Z_1 + Z_1Z_2 + Z_0Z_2\big)$. Its smallest eigenvalue
and the number of computational basis states with that eigenvalue are
\par\smallskip\noindent\begin{tabularx}{\linewidth}{@{}*{2}{L}@{}}
\optlabel{A}\,$-\tfrac12$, six basis states & \optlabel{B}\,$-\tfrac12$, three basis states \\[3pt]
\optlabel{C}\,$+\tfrac32$, two basis states & \optlabel{D}\,$-\tfrac32$, two basis states
\end{tabularx}
\end{qitem}
\qsep

\needspace{7\baselineskip}
\begin{qitem}{10}
For the notebook's 15-node graph ($30$ unit edges), a split $z$ is chosen uniformly at random from all $2^{15}$
bitstrings. The variance of the cut value $C(z)$ is
\par\smallskip\noindent\begin{tabularx}{\linewidth}{@{}*{4}{L}@{}}
\optlabel{A}\,$15$ & \optlabel{B}\,$3.75$ & \optlabel{C}\,$7.5$ & \optlabel{D}\,$30$
\end{tabularx}
\end{qitem}
\qsep

\needspace{7\baselineskip}
\begin{qitem}{11}
\textbf{Statement I:} Because all the operators $Z_iZ_j$ commute with one another,\\
\hspace*{1.5em}$e^{-i\gamma H_C} = \prod_{(i,j)\in E} e^{-i\gamma\frac{w_{ij}}{2} Z_iZ_j}$ holds exactly, with the factors in any order.\\[1mm]
\textbf{Statement II:} Because the operators $X_i$ also commute with one another,\\
\hspace*{1.5em}$e^{-i\beta H_M}\,e^{-i\gamma H_C} = e^{-i(\beta H_M + \gamma H_C)}$ holds for all $\beta, \gamma$.
\begin{itemize}[leftmargin=2.6em, labelsep=0.4em, itemsep=2pt, topsep=3pt]
\item[\optlabel{A}] Both Statement I and Statement II are true
\item[\optlabel{B}] Both Statement I and Statement II are false
\item[\optlabel{C}] Statement I is true, but Statement II is false
\item[\optlabel{D}] Statement I is false, but Statement II is true
\end{itemize}
\end{qitem}
\qsep

\needspace{7\baselineskip}
\begin{qitem}{12}
The expectation value $\langle Z_iZ_j\rangle$ is estimated as the average of the $\pm1$ values $(-1)^{z_i \oplus z_j}$ over
$S$ independent shots. The smallest $S$ that guarantees a standard error of at most $0.01$ \emph{for every state} is
\par\smallskip\noindent\begin{tabularx}{\linewidth}{@{}*{4}{L}@{}}
\optlabel{A}\,$10\,000$ & \optlabel{B}\,$2\,500$ & \optlabel{C}\,$100$ & \optlabel{D}\,$40\,000$
\end{tabularx}
\end{qitem}
\qsep

\needspace{7\baselineskip}
\begin{qitem}{13}
In the previously recorded notebook run, an optimal split is sampled with probability $P_{\text{opt}} = 0.06$ per
shot. Treat the shots as independent and match each quantity (List-I) with its value (List-II).
\par\smallskip\noindent{\renewcommand{\arraystretch}{1.35}%
\begin{tabularx}{\linewidth}{@{}c>{\raggedright\arraybackslash}p{0.33\linewidth}@{\hspace{1.2em}}cL@{}}
\toprule
\multicolumn{2}{@{}l}{\textbf{List-I}} & \multicolumn{2}{l@{}}{\textbf{List-II}} \\ \midrule
\textbf{(P)} & standard deviation of the number of optimal shots among $4000$ & \textbf{(1)} & $75$ \\
\textbf{(Q)} & probability that $10$ shots contain \emph{no} optimal split & \textbf{(2)} & $15.0$ \\
\textbf{(R)} & probability that $10$ shots contain \emph{at least one} optimal split & \textbf{(3)} & $0.539$ \\
\textbf{(S)} & smallest number of shots giving $\ge 99\,\%$ probability of at least one optimal split & \textbf{(4)} & $0.461$ \\
\bottomrule
\end{tabularx}}
\par\smallskip\noindent\begin{tabularx}{\linewidth}{@{}*{2}{L}@{}}
\optlabel{A}\,P$\to$2, Q$\to$3, R$\to$4, S$\to$1 & \optlabel{B}\,P$\to$2, Q$\to$4, R$\to$3, S$\to$1 \\[3pt]
\optlabel{C}\,P$\to$1, Q$\to$4, R$\to$3, S$\to$2 & \optlabel{D}\,P$\to$1, Q$\to$3, R$\to$4, S$\to$2
\end{tabularx}
\end{qitem}
\qsep

\needspace{7\baselineskip}
\begin{qitem}{14}
The notebook's brute-force search reports four optimal bitstrings (cut value $25$):
\[ \texttt{001101110100010},\quad \texttt{001101110110010},\quad \texttt{110010001001101},\quad \texttt{110010001011101}. \]
The number of distinct optimal partitions, and the bitstring describing the same partition as $\texttt{001101110100010}$, are
\par\smallskip\noindent\begin{tabularx}{\linewidth}{@{}*{2}{L}@{}}
\optlabel{A}\,$2$; $\ \texttt{110010001001101}$ & \optlabel{B}\,$2$; $\ \texttt{010001011101100}$ \\[3pt]
\optlabel{C}\,$4$; $\ \texttt{110010001011101}$ & \optlabel{D}\,$2$; $\ \texttt{110010001011101}$
\end{tabularx}
\end{qitem}
\qsep

\needspace{7\baselineskip}
\begin{qitem}{15}
The QUBO function $f(x_0, x_1) = 2x_0x_1 - x_0 - x_1$ with $x_i \in \{0,1\}$ is rewritten with spins via
$x_i = (1 - s_i)/2$, $s_i \in \{+1, -1\}$. The result is
\par\smallskip\noindent\begin{tabularx}{\linewidth}{@{}*{2}{L}@{}}
\optlabel{A}\,$\tfrac12\,s_0s_1 - 1 + \tfrac12(s_0 + s_1)$ & \optlabel{B}\,$\tfrac12\,(s_0s_1 - 1)$ \\[3pt]
\optlabel{C}\,$\tfrac14\,(s_0s_1 - 1)$ & \optlabel{D}\,$\tfrac12\,(1 - s_0s_1)$
\end{tabularx}
\end{qitem}
\qsep

\needspace{7\baselineskip}
\begin{qitem}{16}
Let $f(\theta) = a + b\cos\theta + c\sin\theta$ with constants $a, b, c$. Which expression equals $f'(\theta)$
\emph{exactly} for every $\theta$?
\par\smallskip\noindent\begin{tabularx}{\linewidth}{@{}*{2}{L}@{}}
\optlabel{A}\,$\tfrac12\big[f(\theta + \tfrac\pi4) - f(\theta - \tfrac\pi4)\big]$ & \optlabel{B}\,$f(\theta + \tfrac\pi2) - f(\theta)$ \\[3pt]
\optlabel{C}\,$\tfrac12\big[f(\theta + \tfrac\pi2) - f(\theta - \tfrac\pi2)\big]$ & \optlabel{D}\,$\tfrac1\pi\big[f(\theta + \tfrac\pi2) - f(\theta - \tfrac\pi2)\big]$
\end{tabularx}
\end{qitem}
\qsep

\needspace{20\baselineskip}
\begin{passagebox}{Paragraph for Questions 17 and 18}
The ``house graph'' has the $5$ nodes $0, 1, 2, 3, 4$ and the $6$ unit-weight edges
$(0,1), (1,2), (2,3), (3,0)$ (a square) and $(0,4), (1,4)$ (a roof: node 4 forms a triangle with nodes 0 and 1).
A split is a bitstring $z = z_0z_1z_2z_3z_4$; its cut value $C(z)$ counts the edges whose two end nodes are in different
groups.
\end{passagebox}
\needspace{7\baselineskip}
\begin{qitem}{17}
The maximum cut $C_{\max}$ of the house graph, and the number of bitstrings $z$ that achieve it, are
\par\smallskip\noindent\begin{tabularx}{\linewidth}{@{}*{4}{L}@{}}
\optlabel{A}\,$5$ and $6$ & \optlabel{B}\,$5$ and $2$ & \optlabel{C}\,$6$ and $2$ & \optlabel{D}\,$5$ and $4$
\end{tabularx}
\end{qitem}
\qsep

\needspace{7\baselineskip}
\begin{qitem}{18}
A split of the house graph is drawn uniformly at random. The probability that it is optimal, and the ratio
$\mathbb{E}[C]/C_{\max}$ of the random split, are
\par\smallskip\noindent\begin{tabularx}{\linewidth}{@{}*{4}{L}@{}}
\optlabel{A}\,$\tfrac1{16}$ and $0.6$ & \optlabel{B}\,$\tfrac14$ and $0.6$ & \optlabel{C}\,$\tfrac18$ and $0.5$ & \optlabel{D}\,$\tfrac18$ and $0.6$
\end{tabularx}
\end{qitem}
\qsep

\needspace{7\baselineskip}
\begin{qitem}{19}
For the Petersen graph (10 nodes, 15 unit edges, $C_{\max} = 12$), the smallest eigenvalue of
$H_C = \sum_{(i,j)\in E}\tfrac12 Z_iZ_j$ is
\par\smallskip\noindent\begin{tabularx}{\linewidth}{@{}*{4}{L}@{}}
\optlabel{A}\,$-7.5$ & \optlabel{B}\,$-4.5$ & \optlabel{C}\,$-12$ & \optlabel{D}\,$+4.5$
\end{tabularx}
\end{qitem}
\qsep

\needspace{7\baselineskip}
\begin{qitem}{20}
The notebook's \texttt{local\_search\_maxcut} repeatedly flips single nodes while this strictly increases the cut
(unit weights), starting from random splits.\\[1mm]
\textbf{Assertion (A):} Every split returned by one run of this greedy procedure is an optimal split.\\[1mm]
\textbf{Reason (R):} When no single flip increases the cut, every node has at least half of its incident edges cut.
\begin{itemize}[leftmargin=2.6em, labelsep=0.4em, itemsep=2pt, topsep=3pt]
\item[\optlabel{A}] Both Assertion (A) and Reason (R) are true, and R is the correct explanation of A
\item[\optlabel{B}] Both A and R are true, but R is \emph{not} the correct explanation of A
\item[\optlabel{C}] A is true, but R is false
\item[\optlabel{D}] A is false, but R is true
\end{itemize}
\end{qitem}
\qsep

\newpage

\begin{sectionbanner}SECTION B \enspace$\cdot$\enspace Quantum Mechanics \hfill Q.21--40\end{sectionbanner}

{\small \textbf{Syllabus:} qubits, superposition and the Born rule, Pauli algebra, rotation gates, CNOT conjugation, expectation values, entanglement, global phases, symmetries and state-vector memory.\\ \textbf{Formats:} Single-correct MCQ (Q.21--37); Matching lists (Q.23--27); Assertion--Reason (Q.35); Statement I / II (Q.29--38); Paragraph based (Q.39--40). Each question: $+3$ / $0$ / $-1$.}
\par\medskip

\needspace{7\baselineskip}
\begin{qitem}{21}
The 15-qubit state $\ket{+}^{\otimes 15}$ is measured in the computational basis. The notebook's graph has 4 optimal
bitstrings. The probability of obtaining an optimal bitstring, and the factor by which the previously recorded QAOA value
$P_{\text{opt}} = 0.06$ exceeds it, are
\par\smallskip\noindent\begin{tabularx}{\linewidth}{@{}*{2}{L}@{}}
\optlabel{A}\,$6.10\times10^{-5}$; about $980$ & \optlabel{B}\,$3.05\times10^{-5}$; about $1970$ \\[3pt]
\optlabel{C}\,$2.44\times10^{-4}$; about $245$ & \optlabel{D}\,$1.22\times10^{-4}$; about $490$
\end{tabularx}
\end{qitem}
\qsep

\needspace{7\baselineskip}
\begin{qitem}{22}
The two-qubit operator $(Z\otimes Z)(X\otimes X)$ equals
\par\smallskip\noindent\begin{tabularx}{\linewidth}{@{}*{4}{L}@{}}
\optlabel{A}\,$i\,Y\otimes Y$ & \optlabel{B}\,$-\,Y\otimes Y$ & \optlabel{C}\,$Y\otimes Y$ & \optlabel{D}\,$-i\,Y\otimes Y$
\end{tabularx}
\end{qitem}
\qsep

\needspace{7\baselineskip}
\begin{qitem}{23}
Match each operator (List-I) with its matrix (List-II). Two-qubit matrices use the basis order $\ket{00}, \ket{01},
\ket{10}, \ket{11}$.
\par\smallskip\noindent{\renewcommand{\arraystretch}{1.35}%
\begin{tabularx}{\linewidth}{@{}c>{\raggedright\arraybackslash}p{0.33\linewidth}@{\hspace{1.2em}}cL@{}}
\toprule
\multicolumn{2}{@{}l}{\textbf{List-I}} & \multicolumn{2}{l@{}}{\textbf{List-II}} \\ \midrule
\textbf{(P)} & $R_X(2\beta)$ & \textbf{(1)} & $\operatorname{diag}\big(e^{-i\theta/2}, e^{i\theta/2}, e^{i\theta/2}, e^{-i\theta/2}\big)$ \\
\textbf{(Q)} & $R_Z(\theta)$ & \textbf{(2)} & $\begin{pmatrix}\cos\beta & -i\sin\beta\\ -i\sin\beta & \cos\beta\end{pmatrix}$ \\
\textbf{(R)} & $e^{-i\theta\, Z\otimes Z/2}$ & \textbf{(3)} & $\tfrac{1}{\sqrt2}\begin{pmatrix}1 & 1\\ 1 & -1\end{pmatrix}$ \\
\textbf{(S)} & $H$ & \textbf{(4)} & $\operatorname{diag}\big(e^{-i\theta/2}, e^{i\theta/2}\big)$ \\
\bottomrule
\end{tabularx}}
\par\smallskip\noindent\begin{tabularx}{\linewidth}{@{}*{2}{L}@{}}
\optlabel{A}\,P$\to$3, Q$\to$4, R$\to$1, S$\to$2 & \optlabel{B}\,P$\to$2, Q$\to$1, R$\to$4, S$\to$3 \\[3pt]
\optlabel{C}\,P$\to$4, Q$\to$2, R$\to$1, S$\to$3 & \optlabel{D}\,P$\to$2, Q$\to$4, R$\to$1, S$\to$3
\end{tabularx}
\end{qitem}
\qsep

\needspace{7\baselineskip}
\begin{qitem}{24}
The state $R_Z(\theta)\ket{+}$ is passed through a Hadamard gate and then measured in the computational basis. The
probability of the outcome 0 is
\par\smallskip\noindent\begin{tabularx}{\linewidth}{@{}*{4}{L}@{}}
\optlabel{A}\,$\cos^2(\theta/2)$ & \optlabel{B}\,$\sin^2(\theta/2)$ & \optlabel{C}\,$\cos^2\theta$ & \optlabel{D}\,$\tfrac12$
\end{tabularx}
\end{qitem}
\qsep

\needspace{7\baselineskip}
\begin{qitem}{25}
For one unit edge $(0,1)$ the qubits are in the state $\ket\psi = \tfrac12\ket{00} + \tfrac{\sqrt3}{2}\ket{01}$. The expected
cut value $\langle C\rangle = \big\langle \tfrac12(1 - Z_0Z_1)\big\rangle$ is
\par\smallskip\noindent\begin{tabularx}{\linewidth}{@{}*{4}{L}@{}}
\optlabel{A}\,$\tfrac14$ & \optlabel{B}\,$-\tfrac12$ & \optlabel{C}\,$\tfrac{\sqrt3}{2}$ & \optlabel{D}\,$\tfrac34$
\end{tabularx}
\end{qitem}
\qsep

\needspace{7\baselineskip}
\begin{qitem}{26}
Which of the following operators commutes with $Z_0Z_1$ (on two or more qubits)?
\par\smallskip\noindent\begin{tabularx}{\linewidth}{@{}*{4}{L}@{}}
\optlabel{A}\,$X_0X_1$ & \optlabel{B}\,$Y_0Z_1$ & \optlabel{C}\,$X_0$ & \optlabel{D}\,$X_1$
\end{tabularx}
\end{qitem}
\qsep

\needspace{7\baselineskip}
\begin{qitem}{27}
Conjugation by a CNOT with control qubit 0 and target qubit 1 maps a Pauli operator $P$ to
$\mathrm{CNOT}\,P\,\mathrm{CNOT}$. Match each $P$ (List-I) with its image (List-II).
\par\smallskip\noindent{\renewcommand{\arraystretch}{1.35}%
\begin{tabularx}{\linewidth}{@{}c>{\raggedright\arraybackslash}p{0.33\linewidth}@{\hspace{1.2em}}cL@{}}
\toprule
\multicolumn{2}{@{}l}{\textbf{List-I}} & \multicolumn{2}{l@{}}{\textbf{List-II}} \\ \midrule
\textbf{(P)} & $X\otimes I$ & \textbf{(1)} & $Z\otimes Z$ \\
\textbf{(Q)} & $I\otimes X$ & \textbf{(2)} & $X\otimes X$ \\
\textbf{(R)} & $Z\otimes I$ & \textbf{(3)} & $I\otimes X$ \\
\textbf{(S)} & $I\otimes Z$ & \textbf{(4)} & $Z\otimes I$ \\
\bottomrule
\end{tabularx}}
\par\smallskip\noindent\begin{tabularx}{\linewidth}{@{}*{2}{L}@{}}
\optlabel{A}\,P$\to$3, Q$\to$2, R$\to$1, S$\to$4 & \optlabel{B}\,P$\to$2, Q$\to$3, R$\to$1, S$\to$4 \\[3pt]
\optlabel{C}\,P$\to$3, Q$\to$2, R$\to$4, S$\to$1 & \optlabel{D}\,P$\to$2, Q$\to$3, R$\to$4, S$\to$1
\end{tabularx}
\end{qitem}
\qsep

\needspace{7\baselineskip}
\begin{qitem}{28}
The state $e^{-i\gamma\, Z\otimes Z/2}\ket{+}\ket{+}$ is maximally entangled for which $\gamma \in (0, \pi)$?
\par\smallskip\noindent\begin{tabularx}{\linewidth}{@{}*{2}{L}@{}}
\optlabel{A}\,$\gamma = \pi/2$ & \optlabel{B}\,$\gamma = 3\pi/4$ \\[3pt]
\optlabel{C}\,$\gamma = \pi/4$ & \optlabel{D}\,for no $\gamma$: it is always a product state
\end{tabularx}
\end{qitem}
\qsep

\needspace{7\baselineskip}
\begin{qitem}{29}
\textbf{Statement I:} Replacing $H_C$ by $H_C + c\,I$ (any real constant $c$) changes the QAOA state
$\ket{\boldsymbol\gamma, \boldsymbol\beta}$ only by a global phase.\\[1mm]
\textbf{Statement II:} Consequently, a QAOA circuit that uses $e^{-i\gamma C}$ with $C = \tfrac W2 - H_C$ has the same optimal
angles $(\gamma^*, \beta^*)$ as the notebook's circuit with $e^{-i\gamma H_C}$.
\begin{itemize}[leftmargin=2.6em, labelsep=0.4em, itemsep=2pt, topsep=3pt]
\item[\optlabel{A}] Both Statement I and Statement II are true
\item[\optlabel{B}] Both Statement I and Statement II are false
\item[\optlabel{C}] Statement I is true, but Statement II is false
\item[\optlabel{D}] Statement I is false, but Statement II is true
\end{itemize}
\end{qitem}
\qsep

\needspace{7\baselineskip}
\begin{qitem}{30}
Two qubits are in a product state in which $P(z_0 = 1) = 0.2$ and $P(z_1 = 1) = 0.7$. The value of $\langle Z_0Z_1\rangle$ is
\par\smallskip\noindent\begin{tabularx}{\linewidth}{@{}*{4}{L}@{}}
\optlabel{A}\,$+0.24$ & \optlabel{B}\,$-0.24$ & \optlabel{C}\,$0.14$ & \optlabel{D}\,$0.62$
\end{tabularx}
\end{qitem}
\qsep

\needspace{7\baselineskip}
\begin{qitem}{31}
Conjugating $Z$ by the mixer gate gives $e^{i\beta X}\,Z\,e^{-i\beta X} =$
\par\smallskip\noindent\begin{tabularx}{\linewidth}{@{}*{4}{L}@{}}
\optlabel{A}\,$\cos 2\beta\; Z - \sin 2\beta\; Y$ & \optlabel{B}\,$\cos 2\beta\; Z + \sin 2\beta\; Y$ & \optlabel{C}\,$\cos\beta\; Z + \sin\beta\; Y$ & \optlabel{D}\,$\cos 2\beta\; Z + \sin 2\beta\; X$
\end{tabularx}
\end{qitem}
\qsep

\needspace{7\baselineskip}
\begin{qitem}{32}
A state vector needs $16 \times 2^n$ bytes in fp64 and $8\times 2^n$ bytes in fp32. Ignoring simulator workspace, the
largest $n$ whose state fits on a single 80 GB ($80\times10^9$ bytes) GPU is
\par\smallskip\noindent\begin{tabularx}{\linewidth}{@{}*{2}{L}@{}}
\optlabel{A}\,$33$ (fp64) and $34$ (fp32) & \optlabel{B}\,$36$ (fp64) and $37$ (fp32) \\[3pt]
\optlabel{C}\,$32$ (fp64) and $33$ (fp32) & \optlabel{D}\,$32$ (fp64) and $64$ (fp32)
\end{tabularx}
\end{qitem}
\qsep

\needspace{7\baselineskip}
\begin{qitem}{33}
The gate $R_X(2\beta)$ with $\beta = \pi/8$ is applied to $\ket0$. The probability of measuring 1 is
\par\smallskip\noindent\begin{tabularx}{\linewidth}{@{}*{2}{L}@{}}
\optlabel{A}\,$\sin(\pi/8) \approx 0.383$ & \optlabel{B}\,$\tfrac12$ \\[3pt]
\optlabel{C}\,$\tfrac{2 - \sqrt2}{4} \approx 0.146$ & \optlabel{D}\,$\tfrac{2 + \sqrt2}{4} \approx 0.854$
\end{tabularx}
\end{qitem}
\qsep

\needspace{7\baselineskip}
\begin{qitem}{34}
Which statement about $\ket{+}^{\otimes n}$ and $H_M = \sum_{i} X_i$ is correct?
\begin{itemize}[leftmargin=2.6em, labelsep=0.4em, itemsep=2pt, topsep=3pt]
\item[\optlabel{A}] It is an eigenstate of $H_M$ with eigenvalue $0$
\item[\optlabel{B}] It is the eigenstate of $H_M$ with the \emph{largest} eigenvalue, $+n$
\item[\optlabel{C}] It is not an eigenstate of $H_M$
\item[\optlabel{D}] It is the ground state of $H_M$, with eigenvalue $-n$
\end{itemize}
\end{qitem}
\qsep

\needspace{7\baselineskip}
\begin{qitem}{35}
\textbf{Assertion (A):} A state $\ket\psi$ and its complex conjugate $\ket{\psi^*}$ (every amplitude conjugated) give
identical computational-basis measurement statistics.\\[1mm]
\textbf{Reason (R):} The initial QAOA state $\ket{+}^{\otimes n}$ has only real amplitudes.
\begin{itemize}[leftmargin=2.6em, labelsep=0.4em, itemsep=2pt, topsep=3pt]
\item[\optlabel{A}] Both Assertion (A) and Reason (R) are true, and R is the correct explanation of A
\item[\optlabel{B}] Both A and R are true, but R is \emph{not} the correct explanation of A
\item[\optlabel{C}] A is true, but R is false
\item[\optlabel{D}] A is false, but R is true
\end{itemize}
\end{qitem}
\qsep

\needspace{7\baselineskip}
\begin{qitem}{36}
For the triangle $K_3$ (unit weights) a state has $\langle Z_0Z_1\rangle = -0.6$, $\langle Z_1Z_2\rangle = -0.2$ and
$\langle Z_0Z_2\rangle = 0.4$. Its expected cut $\langle C\rangle$ is
\par\smallskip\noindent\begin{tabularx}{\linewidth}{@{}*{4}{L}@{}}
\optlabel{A}\,$-0.2$ & \optlabel{B}\,$1.9$ & \optlabel{C}\,$1.3$ & \optlabel{D}\,$1.7$
\end{tabularx}
\end{qitem}
\qsep

\needspace{7\baselineskip}
\begin{qitem}{37}
On $n$ qubits, $e^{-i(\pi/2) H_M}$ with $H_M = \sum_i X_i$ equals ..., and therefore (since $C(z) = C(\bar z)$) $F_p$ is
periodic in every single $\beta_l$ with period ...
\par\smallskip\noindent\begin{tabularx}{\linewidth}{@{}*{2}{L}@{}}
\optlabel{A}\,$(-i)^n X^{\otimes n}$; $\ \pi/2$ & \optlabel{B}\,$(-i)^n X^{\otimes n}$; $\ \pi$ \\[3pt]
\optlabel{C}\,$I$; $\ \pi/2$ & \optlabel{D}\,$X^{\otimes n}$; $\ 2\pi$
\end{tabularx}
\end{qitem}
\qsep

\needspace{7\baselineskip}
\begin{qitem}{38}
\textbf{Statement I:} $Z_iZ_j$ has eigenvalue $-1$ exactly on the computational basis states with $z_i \neq z_j$.\\[1mm]
\textbf{Statement II:} $\langle Z_iZ_j\rangle$ can be estimated from the same computational-basis shots that are used to sample
candidate splits.
\begin{itemize}[leftmargin=2.6em, labelsep=0.4em, itemsep=2pt, topsep=3pt]
\item[\optlabel{A}] Both Statement I and Statement II are true
\item[\optlabel{B}] Both Statement I and Statement II are false
\item[\optlabel{C}] Statement I is true, but Statement II is false
\item[\optlabel{D}] Statement I is false, but Statement II is true
\end{itemize}
\end{qitem}
\qsep

\needspace{20\baselineskip}
\begin{passagebox}{Paragraph for Questions 39 and 40}
A single edge $(0,1)$ with weight 1 is treated with one QAOA layer in the notebook's convention. Starting from
$\ket{00}$ the gates are: $H$ on both qubits; CNOT (control 0, target 1); $R_Z(\theta)$ on qubit 1 with $\theta = \gamma w = \gamma$;
CNOT (control 0, target 1); $R_X(2\beta)$ on both qubits. The angles are $\gamma = \pi/2$ and $\beta = \pi/8$.
Amplitudes are listed in the order $\ket{00}, \ket{01}, \ket{10}, \ket{11}$ (qubit 0 written first).
\end{passagebox}
\needspace{7\baselineskip}
\begin{qitem}{39}
Immediately after the second CNOT (before the mixer), the phase of the amplitude of $\ket{01}$ relative to the amplitude
of $\ket{00}$ is
\par\smallskip\noindent\begin{tabularx}{\linewidth}{@{}*{4}{L}@{}}
\optlabel{A}\,$45^\circ$ & \optlabel{B}\,$180^\circ$ & \optlabel{C}\,$90^\circ$ & \optlabel{D}\,$0^\circ$
\end{tabularx}
\end{qitem}
\qsep

\needspace{7\baselineskip}
\begin{qitem}{40}
After the mixer, the probability of measuring a \emph{cut} bitstring ($01$ or $10$) is
\par\smallskip\noindent\begin{tabularx}{\linewidth}{@{}*{4}{L}@{}}
\optlabel{A}\,$1$ & \optlabel{B}\,$\tfrac34$ & \optlabel{C}\,$0$ & \optlabel{D}\,$\tfrac12$
\end{tabularx}
\end{qitem}
\qsep

\newpage

\begin{sectionbanner}SECTION C \enspace$\cdot$\enspace QAOA Basics \hfill Q.41--60\end{sectionbanner}

{\small \textbf{Syllabus:} cost and mixer Hamiltonians, initial state, QAOA circuit and gate decomposition, parameters and depth, the $p = 1$ analysis, objective function, sign conventions and the notebook's code.\\ \textbf{Formats:} Single-correct MCQ (Q.41--57); Matching lists (Q.46--55); Assertion--Reason (Q.53); Statement I / II (Q.49--58); Paragraph based (Q.59--60). Each question: $+3$ / $0$ / $-1$.}
\par\medskip

\needspace{7\baselineskip}
\begin{qitem}{41}
For \texttt{edges = [(0, 1, 2.0), (1, 2, 1.0)]} the notebook's function \texttt{cost\_hamiltonian()} returns the operator
\par\smallskip\noindent\begin{tabularx}{\linewidth}{@{}*{2}{L}@{}}
\optlabel{A}\,$1.0\,Z_0Z_1 + 0.5\,Z_1Z_2$ & \optlabel{B}\,$-1.0\,Z_0Z_1 - 0.5\,Z_1Z_2$ \\[3pt]
\optlabel{C}\,$2.0\,Z_0Z_1 + 1.0\,Z_1Z_2$ & \optlabel{D}\,$1.5\,I - 1.0\,Z_0Z_1 - 0.5\,Z_1Z_2$
\end{tabularx}
\end{qitem}
\qsep

\needspace{7\baselineskip}
\begin{qitem}{42}
In the notebook's convention, if all cost angles are zero ($\gamma_1 = \dots = \gamma_p = 0$), the expected cut
$F_p(\boldsymbol\gamma, \boldsymbol\beta)$ equals
\par\smallskip\noindent\begin{tabularx}{\linewidth}{@{}*{2}{L}@{}}
\optlabel{A}\,$W/2 + \tfrac12\sum_l \sin 4\beta_l$ & \optlabel{B}\,$C_{\max}$ \\[3pt]
\optlabel{C}\,$W/2$ for every choice of $\boldsymbol\beta$ & \optlabel{D}\,$0$
\end{tabularx}
\end{qitem}
\qsep

\needspace{7\baselineskip}
\begin{qitem}{43}
In the notebook's kernel each edge $(i, j)$ with weight $w$ is implemented as
\texttt{x.ctrl(q[i], q[j]); rz(gamma * w, q[j]); x.ctrl(q[i], q[j])}. With $R_Z(\theta) = e^{-i\theta Z/2}$, this sequence equals
\par\smallskip\noindent\begin{tabularx}{\linewidth}{@{}*{4}{L}@{}}
\optlabel{A}\,$e^{+i\gamma\,(w/2)\,Z_iZ_j}$ & \optlabel{B}\,$e^{-i\gamma\, w\, Z_iZ_j}$ & \optlabel{C}\,$e^{-i\gamma\,(w/2)\,Z_j}$ & \optlabel{D}\,$e^{-i\gamma\,(w/2)\,Z_iZ_j}$
\end{tabularx}
\end{qitem}
\qsep

\needspace{7\baselineskip}
\begin{qitem}{44}
Why does the notebook apply \texttt{rx(2.0 * beta, q)} (and not \texttt{rx(beta, q)}) in the mixer layer?
\begin{itemize}[leftmargin=2.6em, labelsep=0.4em, itemsep=2pt, topsep=3pt]
\item[\optlabel{A}] Because the mixer must be applied twice per layer
\item[\optlabel{B}] Because the mixer Hamiltonian is $H_M = 2\sum_i X_i$
\item[\optlabel{C}] Because $\beta$ is measured in units of $\pi/2$
\item[\optlabel{D}] Because $R_X(\theta) = e^{-i\theta X/2}$, so $R_X(2\beta) = e^{-i\beta X}$, the factor of $e^{-i\beta H_M}$ for one qubit
\end{itemize}
\end{qitem}
\qsep

\needspace{7\baselineskip}
\begin{qitem}{45}
Using the notebook's gate structure ($H$ on every qubit, then per layer 3 gates per edge and one $R_X$ per qubit), a
graph with $n = 20$ nodes and $|E| = 57$ edges at depth $p = 3$ needs how many gates (excluding measurements), and how many of
them are CNOTs?
\par\smallskip\noindent\begin{tabularx}{\linewidth}{@{}*{4}{L}@{}}
\optlabel{A}\,$573$ and $342$ & \optlabel{B}\,$593$ and $342$ & \optlabel{C}\,$593$ and $171$ & \optlabel{D}\,$402$ and $228$
\end{tabularx}
\end{qitem}
\qsep

\needspace{7\baselineskip}
\begin{qitem}{46}
For the notebook's default run ($n = 15$ nodes, $|E| = 30$ edges, $p = 2$), match each quantity (List-I) with its value
(List-II).
\par\smallskip\noindent{\renewcommand{\arraystretch}{1.35}%
\begin{tabularx}{\linewidth}{@{}c>{\raggedright\arraybackslash}p{0.33\linewidth}@{\hspace{1.2em}}cL@{}}
\toprule
\multicolumn{2}{@{}l}{\textbf{List-I}} & \multicolumn{2}{l@{}}{\textbf{List-II}} \\ \midrule
\textbf{(P)} & number of variational parameters & \textbf{(1)} & $105$ \\
\textbf{(Q)} & number of $R_Z$ gates & \textbf{(2)} & $120$ \\
\textbf{(R)} & number of single-qubit gates ($H$, $R_Z$ and $R_X$ together) & \textbf{(3)} & $4$ \\
\textbf{(S)} & number of CNOT (\texttt{cx}) gates & \textbf{(4)} & $60$ \\
\bottomrule
\end{tabularx}}
\par\smallskip\noindent\begin{tabularx}{\linewidth}{@{}*{2}{L}@{}}
\optlabel{A}\,P$\to$4, Q$\to$3, R$\to$1, S$\to$2 & \optlabel{B}\,P$\to$3, Q$\to$4, R$\to$2, S$\to$1 \\[3pt]
\optlabel{C}\,P$\to$3, Q$\to$4, R$\to$1, S$\to$2 & \optlabel{D}\,P$\to$3, Q$\to$2, R$\to$1, S$\to$4
\end{tabularx}
\end{qitem}
\qsep

\needspace{7\baselineskip}
\begin{qitem}{47}
For a single unit edge and $p = 1$ in the notebook's convention, the expected cut is
$F_1(\gamma, \beta) = \tfrac12 - \tfrac12\sin4\beta\,\sin\gamma$. Its maximum value, and one point where it is reached, are
\par\smallskip\noindent\begin{tabularx}{\linewidth}{@{}*{2}{L}@{}}
\optlabel{A}\,$\tfrac34$, at $(\gamma, \beta) = (\pi/2,\, -\pi/8)$ & \optlabel{B}\,$\tfrac12$, reached for all angles \\[3pt]
\optlabel{C}\,$1$, at $(\gamma, \beta) = (\pi/2,\, \pi/8)$ & \optlabel{D}\,$1$, at $(\gamma, \beta) = (\pi/2,\, -\pi/8)$
\end{tabularx}
\end{qitem}
\qsep

\needspace{7\baselineskip}
\begin{qitem}{48}
With optimally chosen angles, the best approximation ratio of $p = 1$ QAOA on the ring (cycle) $C_8$ is
\par\smallskip\noindent\begin{tabularx}{\linewidth}{@{}*{4}{L}@{}}
\optlabel{A}\,$1$ & \optlabel{B}\,$3/4$ & \optlabel{C}\,$5/6$ & \optlabel{D}\,$0.6924$
\end{tabularx}
\end{qitem}
\qsep

\needspace{7\baselineskip}
\begin{qitem}{49}
\textbf{Statement I:} For $p = 1$ with optimal angles, QAOA achieves $\langle C\rangle \ge 0.6924\,C_{\max}$ on every 3-regular graph.\\[1mm]
\textbf{Statement II:} For every graph, the optimal value $F_p^* = \max_{\boldsymbol\gamma,\boldsymbol\beta}F_p$ satisfies $F_{p+1}^* \ge F_p^*$.
\begin{itemize}[leftmargin=2.6em, labelsep=0.4em, itemsep=2pt, topsep=3pt]
\item[\optlabel{A}] Both Statement I and Statement II are true
\item[\optlabel{B}] Both Statement I and Statement II are false
\item[\optlabel{C}] Statement I is true, but Statement II is false
\item[\optlabel{D}] Statement I is false, but Statement II is true
\end{itemize}
\end{qitem}
\qsep

\needspace{7\baselineskip}
\begin{qitem}{50}
Which statement about the limit of large depth $p$ (angles optimised for each $p$) is correct?
\begin{itemize}[leftmargin=2.6em, labelsep=0.4em, itemsep=2pt, topsep=3pt]
\item[\optlabel{A}] $F_p^*/C_{\max} \to 0.878$, the Goemans--Williamson ratio
\item[\optlabel{B}] $F_p^* = C_{\max}$ already at $p = 1$ for every graph
\item[\optlabel{C}] $F_p^* \to W/2$, because deep circuits randomise the state
\item[\optlabel{D}] $F_p^* \to C_{\max}$ as $p \to \infty$, because a Trotterised adiabatic path is one admissible choice of angles
\end{itemize}
\end{qitem}
\qsep

\needspace{7\baselineskip}
\begin{qitem}{51}
In the notebook the angle vector is $\texttt{params} = [\gamma_1, \dots, \gamma_p, \beta_1, \dots, \beta_p]$; the kernel reads
\texttt{gamma = params[layer]} and \texttt{beta = params[layers + layer]}. For $p = 3$ and
\texttt{params = [a, b, c, d, e, f]}, the mixer angle $\beta_2$ of the second layer is
\par\smallskip\noindent\begin{tabularx}{\linewidth}{@{}*{4}{L}@{}}
\optlabel{A}\,$d$ & \optlabel{B}\,$b$ & \optlabel{C}\,$e$ & \optlabel{D}\,$f$
\end{tabularx}
\end{qitem}
\qsep

\needspace{7\baselineskip}
\begin{qitem}{52}
With \texttt{INITIAL\_PARAMETERS = "ramp"} the notebook starts from \texttt{np.linspace(0.2, 0.6, p)} for the $\gamma$'s and
\texttt{np.linspace(0.6, 0.2, p)} for the $\beta$'s. For $p = 5$, the initial values of $\gamma_3$ and $\beta_2$ are
\par\smallskip\noindent\begin{tabularx}{\linewidth}{@{}*{4}{L}@{}}
\optlabel{A}\,$0.3$ and $0.5$ & \optlabel{B}\,$0.36$ and $0.52$ & \optlabel{C}\,$0.4$ and $0.3$ & \optlabel{D}\,$0.4$ and $0.5$
\end{tabularx}
\end{qitem}
\qsep

\needspace{7\baselineskip}
\begin{qitem}{53}
\textbf{Assertion (A):} The notebook's objective function returns $\langle H_C\rangle$ and SciPy \emph{minimises} it; this is
equivalent to \emph{maximising} the expected cut.\\[1mm]
\textbf{Reason (R):} $\langle C\rangle = W/2 - \langle H_C\rangle$, and $W/2$ does not depend on the angles.
\begin{itemize}[leftmargin=2.6em, labelsep=0.4em, itemsep=2pt, topsep=3pt]
\item[\optlabel{A}] Both Assertion (A) and Reason (R) are true, and R is the correct explanation of A
\item[\optlabel{B}] Both A and R are true, but R is \emph{not} the correct explanation of A
\item[\optlabel{C}] A is true, but R is false
\item[\optlabel{D}] A is false, but R is true
\end{itemize}
\end{qitem}
\qsep

\needspace{7\baselineskip}
\begin{qitem}{54}
At the notebook's default starting point $\boldsymbol\gamma = (0.2, 0.6)$, $\boldsymbol\beta = (0.6, 0.2)$, an exact
state-vector calculation for its 15-node graph gives $\langle C\rangle \approx 9.18$, \emph{below} the random baseline $W/2 = 15$.
The best explanation is
\begin{itemize}[leftmargin=2.6em, labelsep=0.4em, itemsep=2pt, topsep=3pt]
\item[\optlabel{A}] In the convention $e^{-i\beta H_M}e^{-i\gamma H_C}$, the start state $\ket+^{\otimes n}$ is the \emph{highest} eigenstate of $H_M$, so an annealing-like ramp with $\gamma, \beta > 0$ heads for the highest eigenstate of $H_C$, the \emph{smallest} cut; heading for the largest cut needs $\gamma$ and $\beta$ of opposite sign
\item[\optlabel{B}] The value is meaningless: \texttt{cudaq.observe} returns random numbers before the optimizer has converged
\item[\optlabel{C}] Every $p = 2$ starting point gives $\langle C\rangle < W/2$; only $p \ge 3$ can beat random guessing
\item[\optlabel{D}] It indicates a bug, because a QAOA state can never have $\langle C\rangle < W/2$
\end{itemize}
\end{qitem}
\qsep

\needspace{7\baselineskip}
\begin{qitem}{55}
Match each part of the QAOA algorithm (List-I) with the code that implements it in the notebook (List-II).
\par\smallskip\noindent{\renewcommand{\arraystretch}{1.35}%
\begin{tabularx}{\linewidth}{@{}c>{\raggedright\arraybackslash}p{0.33\linewidth}@{\hspace{1.2em}}cL@{}}
\toprule
\multicolumn{2}{@{}l}{\textbf{List-I}} & \multicolumn{2}{l@{}}{\textbf{List-II}} \\ \midrule
\textbf{(P)} & initial state $\ket+^{\otimes n}$ & \textbf{(1)} & \texttt{rx(2.0 * beta, qubits[q])} \\
\textbf{(Q)} & cost unitary for one edge & \textbf{(2)} & \texttt{cudaq.observe(kernels["state"], H\_C, ...).expectation()} \\
\textbf{(R)} & mixer unitary for one qubit & \textbf{(3)} & \texttt{h(qubits)} \\
\textbf{(S)} & objective $\langle H_C\rangle$ & \textbf{(4)} & \texttt{x.ctrl(...); rz(gamma * weights[e], ...); x.ctrl(...)} \\
\bottomrule
\end{tabularx}}
\par\smallskip\noindent\begin{tabularx}{\linewidth}{@{}*{2}{L}@{}}
\optlabel{A}\,P$\to$3, Q$\to$1, R$\to$4, S$\to$2 & \optlabel{B}\,P$\to$3, Q$\to$4, R$\to$2, S$\to$1 \\[3pt]
\optlabel{C}\,P$\to$3, Q$\to$4, R$\to$1, S$\to$2 & \optlabel{D}\,P$\to$4, Q$\to$3, R$\to$1, S$\to$2
\end{tabularx}
\end{qitem}
\qsep

\needspace{7\baselineskip}
\begin{qitem}{56}
For the triangle $K_3$ (unit weights, $C_{\max} = 2$), the best value of $F_1$ over all angles is
\par\smallskip\noindent\begin{tabularx}{\linewidth}{@{}*{2}{L}@{}}
\optlabel{A}\,$\sqrt3$, i.e.\ $r \approx 0.866$ & \optlabel{B}\,$2$, i.e.\ $r = 1$ \\[3pt]
\optlabel{C}\,$1.75$, i.e.\ $r = 0.875$ & \optlabel{D}\,$1.5$, i.e.\ $r = 0.75$
\end{tabularx}
\end{qitem}
\qsep

\needspace{7\baselineskip}
\begin{qitem}{57}
The previously recorded run reports: approximation ratio $r = \langle C\rangle/C_{\max} = 0.8213$, best sampled cut $25$
($= C_{\max}$), and $P(\text{optimal split}) = 0.06$ with $4000$ shots. Which interpretation is correct?
\begin{itemize}[leftmargin=2.6em, labelsep=0.4em, itemsep=2pt, topsep=3pt]
\item[\optlabel{A}] $82\,\%$ of the shots returned an optimal split
\item[\optlabel{B}] $r$ measures the \emph{average} cut of the distribution; the best of $4000$ samples can still be optimal, and here $240$ shots were optimal
\item[\optlabel{C}] The best sampled cut can never exceed $\langle C\rangle = 20.53$, so the value $25$ indicates an error
\item[\optlabel{D}] The best sampled split has cut $0.8213 \times 25 \approx 20.5$
\end{itemize}
\end{qitem}
\qsep

\needspace{7\baselineskip}
\begin{qitem}{58}
\textbf{Statement I:} The notebook's final sampling run uses the angles of the best evaluation stored in
\texttt{history\_cut}, which need not be the optimizer's final point \texttt{opt.x}.\\[1mm]
\textbf{Statement II:} With COBYLA and a budget of 80 evaluations, the notebook is guaranteed to find the global maximum of $F_2$.
\begin{itemize}[leftmargin=2.6em, labelsep=0.4em, itemsep=2pt, topsep=3pt]
\item[\optlabel{A}] Both Statement I and Statement II are true
\item[\optlabel{B}] Both Statement I and Statement II are false
\item[\optlabel{C}] Statement I is true, but Statement II is false
\item[\optlabel{D}] Statement I is false, but Statement II is true
\end{itemize}
\end{qitem}
\qsep

\needspace{20\baselineskip}
\begin{passagebox}{Paragraph for Questions 59 and 60}
QAOA with $p = 1$ is applied to the path graph $0 - 1 - 2$ (edges $(0,1)$ and $(1,2)$, unit weights,
$C_{\max} = 2$) in the notebook's convention, with $\gamma = \pi/2$ and $\beta = -\pi/8$. For a triangle-free graph and $p = 1$
the edge correlator in this convention is
\[ \langle Z_uZ_v\rangle = \tfrac12\,\sin4\beta\,\sin\gamma\,\big(\cos^{d_u - 1}\gamma + \cos^{d_v - 1}\gamma\big), \]
where $d_u$, $d_v$ are the degrees of the two end nodes.
\end{passagebox}
\needspace{7\baselineskip}
\begin{qitem}{59}
The value of $\langle Z_0Z_1\rangle$ is
\par\smallskip\noindent\begin{tabularx}{\linewidth}{@{}*{4}{L}@{}}
\optlabel{A}\,$-\tfrac12$ & \optlabel{B}\,$-1$ & \optlabel{C}\,$+\tfrac12$ & \optlabel{D}\,$0$
\end{tabularx}
\end{qitem}
\qsep

\needspace{7\baselineskip}
\begin{qitem}{60}
The expected cut $\langle C\rangle$ and the probability of sampling an optimal split are
\par\smallskip\noindent\begin{tabularx}{\linewidth}{@{}*{4}{L}@{}}
\optlabel{A}\,$1.5$ and $0.3125$ & \optlabel{B}\,$1.5$ and $0.625$ & \optlabel{C}\,$1.25$ and $0.625$ & \optlabel{D}\,$2$ and $1$
\end{tabularx}
\end{qitem}
\qsep

\newpage

\begin{sectionbanner}SECTION D \enspace$\cdot$\enspace Advanced QAOA and Implementation \hfill Q.61--80\end{sectionbanner}

{\small \textbf{Syllabus:} light cones, gradients and classical optimizers, shot noise, parameter strategies, simulation backends and memory, noise, QAOA variants, limitations, benchmarking and interpretation of recorded results.\\ \textbf{Formats:} Single-correct MCQ (Q.61--78); Matching lists (Q.68--72); Assertion--Reason (Q.70); Statement I / II (Q.74); Paragraph based (Q.79--80). Each question: $+3$ / $0$ / $-1$.}
\par\medskip

\needspace{7\baselineskip}
\begin{qitem}{61}
On a ring graph, the expectation $\langle Z_uZ_v\rangle$ for one edge $(u, v)$ of a depth-$p$ QAOA state depends on at most how
many qubits when $p = 3$?
\par\smallskip\noindent\begin{tabularx}{\linewidth}{@{}*{4}{L}@{}}
\optlabel{A}\,$8$ & \optlabel{B}\,all $n$ qubits & \optlabel{C}\,$7$ & \optlabel{D}\,$6$
\end{tabularx}
\end{qitem}
\qsep

\needspace{7\baselineskip}
\begin{qitem}{62}
For even rings $C_n$ with $n \ge 2p + 2$, numerical optimisation of the angles gives the best $p = 2$ approximation ratio
\par\smallskip\noindent\begin{tabularx}{\linewidth}{@{}*{4}{L}@{}}
\optlabel{A}\,$3/4$ & \optlabel{B}\,$5/6$ & \optlabel{C}\,$7/8$ & \optlabel{D}\,$1$
\end{tabularx}
\end{qitem}
\qsep

\needspace{7\baselineskip}
\begin{qitem}{63}
For the default circuit ($n = 15$, $|E| = 30$, $p = 2$), compare two ways of computing the gradient of $F_2$ with respect to
all $2p$ angles: (i) the exact parameter-shift rule applied to every rotation gate separately (each $R_Z$ and $R_X$ gate is shifted
by $\pm\pi/2$ on its own), and (ii) a forward finite-difference gradient as used by L-BFGS-B. The numbers of circuit evaluations are
\par\smallskip\noindent\begin{tabularx}{\linewidth}{@{}*{4}{L}@{}}
\optlabel{A}\,$8$ and $5$ & \optlabel{B}\,$180$ and $4$ & \optlabel{C}\,$90$ and $5$ & \optlabel{D}\,$180$ and $5$
\end{tabularx}
\end{qitem}
\qsep

\needspace{7\baselineskip}
\begin{qitem}{64}
The notebook optimises the angles with SciPy's \texttt{minimize(..., method="COBYLA")} and the options
\texttt{\{"maxiter": 80, "rhobeg": 0.2\}}. The option \texttt{rhobeg = 0.2} sets
\begin{itemize}[leftmargin=2.6em, labelsep=0.4em, itemsep=2pt, topsep=3pt]
\item[\optlabel{A}] a learning rate multiplying the gradient
\item[\optlabel{B}] the maximum number of iterations divided by 400
\item[\optlabel{C}] the tolerance at which COBYLA declares convergence
\item[\optlabel{D}] the initial size of the trust region, i.e.\ the first steps change the angles by about $0.2$ rad
\end{itemize}
\end{qitem}
\qsep

\needspace{7\baselineskip}
\begin{qitem}{65}
With \texttt{OBJECTIVE\_SHOTS = 1000}, $\langle H_C\rangle$ is estimated from 1000 shots. If the cut values of the sampled
bitstrings have a standard deviation of $2.5$, the standard error of the estimated objective is about
\par\smallskip\noindent\begin{tabularx}{\linewidth}{@{}*{4}{L}@{}}
\optlabel{A}\,$0.0063$ & \optlabel{B}\,$0.0025$ & \optlabel{C}\,$0.079$ & \optlabel{D}\,$2.5$
\end{tabularx}
\end{qitem}
\qsep

\needspace{7\baselineskip}
\begin{qitem}{66}
For $n = 22$ (\texttt{BRUTE\_FORCE\_MAX\_NODES}), the notebook's \texttt{brute\_force\_maxcut} creates the arrays
\texttt{index = np.arange(2**N, dtype=np.int64)} and \texttt{values = np.zeros(2**N)}. Together they occupy about
\par\smallskip\noindent\begin{tabularx}{\linewidth}{@{}*{4}{L}@{}}
\optlabel{A}\,$34$ MB & \optlabel{B}\,$67$ MB & \optlabel{C}\,$0.5$ MB & \optlabel{D}\,$4.2$ MB
\end{tabularx}
\end{qitem}
\qsep

\needspace{7\baselineskip}
\begin{qitem}{67}
A 36-qubit QAOA state vector is simulated in fp64. If it is split evenly over a power-of-two number of 80 GB GPUs (as the
multi-GPU state-vector target does), ignoring workspace, the minimum number of GPUs is
\par\smallskip\noindent\begin{tabularx}{\linewidth}{@{}*{4}{L}@{}}
\optlabel{A}\,$8$ & \optlabel{B}\,$14$ & \optlabel{C}\,$32$ & \optlabel{D}\,$16$
\end{tabularx}
\end{qitem}
\qsep

\needspace{7\baselineskip}
\begin{qitem}{68}
Match each setting of \texttt{BACKEND} and \texttt{SIMULATION\_METHOD} in \texttt{QAOA.ipynb} (List-I) with what the notebook's
\texttt{resolve\_backend} does (List-II).
\par\smallskip\noindent{\renewcommand{\arraystretch}{1.35}%
\begin{tabularx}{\linewidth}{@{}c>{\raggedright\arraybackslash}p{0.33\linewidth}@{\hspace{1.2em}}cL@{}}
\toprule
\multicolumn{2}{@{}l}{\textbf{List-I}} & \multicolumn{2}{l@{}}{\textbf{List-II}} \\ \midrule
\textbf{(P)} & \texttt{"cpu"} + \texttt{"statevector"} & \textbf{(1)} & target \texttt{tensornet-mps} \\
\textbf{(Q)} & \texttt{"cpu"} + \texttt{"density\_matrix"} & \textbf{(2)} & stops with an error: the MPS simulator needs an NVIDIA GPU \\
\textbf{(R)} & \texttt{"gpu"} + \texttt{"mps"} & \textbf{(3)} & target \texttt{qpp-cpu} \\
\textbf{(S)} & \texttt{"cpu"} + \texttt{"mps"} & \textbf{(4)} & target \texttt{density-matrix-cpu} \\
\bottomrule
\end{tabularx}}
\par\smallskip\noindent\begin{tabularx}{\linewidth}{@{}*{2}{L}@{}}
\optlabel{A}\,P$\to$4, Q$\to$3, R$\to$1, S$\to$2 & \optlabel{B}\,P$\to$3, Q$\to$1, R$\to$4, S$\to$2 \\[3pt]
\optlabel{C}\,P$\to$3, Q$\to$4, R$\to$2, S$\to$1 & \optlabel{D}\,P$\to$3, Q$\to$4, R$\to$1, S$\to$2
\end{tabularx}
\end{qitem}
\qsep

\needspace{7\baselineskip}
\begin{qitem}{69}
For which instance is a matrix-product-state (MPS) simulation of a depth-$p$ QAOA circuit most likely to be efficient (small
bond dimension), with the qubits ordered along a line $0, 1, \dots, n-1$?
\begin{itemize}[leftmargin=2.6em, labelsep=0.4em, itemsep=2pt, topsep=3pt]
\item[\optlabel{A}] Any graph, because QAOA states are always product states before measurement
\item[\optlabel{B}] A complete graph $K_n$: its symmetry makes the MPS exact with bond dimension 1
\item[\optlabel{C}] A ring $C_n$ with small $p$: every cut of the line is crossed by only 2 edges
\item[\optlabel{D}] A dense random graph with edge probability $0.5$: more edges average out the entanglement
\end{itemize}
\end{qitem}
\qsep

\needspace{7\baselineskip}
\begin{qitem}{70}
\textbf{Assertion (A):} When \texttt{BACKEND = "hardware"}, the notebook estimates the objective from shots even if
\texttt{OBJECTIVE\_SHOTS = None}.\\[1mm]
\textbf{Reason (R):} Hardware (and hardware-emulation) targets only return measurement samples; they cannot give a state vector or an
exact expectation value.
\begin{itemize}[leftmargin=2.6em, labelsep=0.4em, itemsep=2pt, topsep=3pt]
\item[\optlabel{A}] Both Assertion (A) and Reason (R) are true, and R is the correct explanation of A
\item[\optlabel{B}] Both A and R are true, but R is \emph{not} the correct explanation of A
\item[\optlabel{C}] A is true, but R is false
\item[\optlabel{D}] A is false, but R is true
\end{itemize}
\end{qitem}
\qsep

\needspace{7\baselineskip}
\begin{qitem}{71}
A simple global-depolarising noise model treats the noisy QAOA output as a mixture: with probability $F = (1 - \varepsilon)^{N_{\text{CNOT}}}$
the ideal state, otherwise a uniformly random bitstring. For the default circuit ($N_{\text{CNOT}} = 120$, ideal $\langle C\rangle = 20.53$,
$W/2 = 15$) and $\varepsilon = 0.005$ per CNOT, the predicted noisy $\langle C\rangle$ is about
\par\smallskip\noindent\begin{tabularx}{\linewidth}{@{}*{4}{L}@{}}
\optlabel{A}\,$17.2$ & \optlabel{B}\,$15.0$ & \optlabel{C}\,$18.0$ & \optlabel{D}\,$11.3$
\end{tabularx}
\end{qitem}
\qsep

\needspace{7\baselineskip}
\begin{qitem}{72}
Match each QAOA variant (List-I) with its key idea (List-II).
\par\smallskip\noindent{\renewcommand{\arraystretch}{1.35}%
\begin{tabularx}{\linewidth}{@{}c>{\raggedright\arraybackslash}p{0.33\linewidth}@{\hspace{1.2em}}cL@{}}
\toprule
\multicolumn{2}{@{}l}{\textbf{List-I}} & \multicolumn{2}{l@{}}{\textbf{List-II}} \\ \midrule
\textbf{(P)} & warm-start QAOA & \textbf{(1)} & problem-specific mixers (e.g.\ XY mixers) that keep the state inside the feasible subspace of a constrained problem \\
\textbf{(Q)} & ADAPT-QAOA & \textbf{(2)} & use measured correlations $\langle Z_iZ_j\rangle$ to fix and eliminate variables, then solve the smaller instance \\
\textbf{(R)} & Quantum Alternating Operator Ansatz (Hadfield et al.) & \textbf{(3)} & start from a state biased towards a classical (e.g.\ relaxed/SDP) solution instead of $\ket+^{\otimes n}$ \\
\textbf{(S)} & recursive QAOA (RQAOA) & \textbf{(4)} & grow the mixer layer by layer from an operator pool, choosing the operator with the largest energy gradient \\
\bottomrule
\end{tabularx}}
\par\smallskip\noindent\begin{tabularx}{\linewidth}{@{}*{2}{L}@{}}
\optlabel{A}\,P$\to$3, Q$\to$1, R$\to$4, S$\to$2 & \optlabel{B}\,P$\to$3, Q$\to$4, R$\to$2, S$\to$1 \\[3pt]
\optlabel{C}\,P$\to$3, Q$\to$4, R$\to$1, S$\to$2 & \optlabel{D}\,P$\to$4, Q$\to$3, R$\to$1, S$\to$2
\end{tabularx}
\end{qitem}
\qsep

\needspace{7\baselineskip}
\begin{qitem}{73}
The INTERP initialisation (Zhou et al.\ 2020) builds depth-$(p+1)$ starting angles from optimised depth-$p$ angles by
$\gamma_i^{(p+1)} = \tfrac{i-1}{p}\,\gamma_{i-1}^{(p)} + \tfrac{p-i+1}{p}\,\gamma_i^{(p)}$ for $i = 1, \dots, p+1$ (with $\gamma_0^{(p)} = \gamma_{p+1}^{(p)} = 0$).
From $\boldsymbol\gamma^{(2)} = (0.4, 0.8)$ it gives $\boldsymbol\gamma^{(3)} =$
\par\smallskip\noindent\begin{tabularx}{\linewidth}{@{}*{2}{L}@{}}
\optlabel{A}\,$(0.4,\ 0.6,\ 0.8)$ & \optlabel{B}\,$(0.27,\ 0.53,\ 0.8)$ \\[3pt]
\optlabel{C}\,$(0.4,\ 0.8,\ 0.8)$ & \optlabel{D}\,$(0.4,\ 0.8,\ 0)$
\end{tabularx}
\end{qitem}
\qsep

\needspace{7\baselineskip}
\begin{qitem}{74}
For graphs with more nodes than \texttt{BRUTE\_FORCE\_MAX\_NODES}, the notebook replaces the exact optimum by the best cut found
by \texttt{local\_search\_maxcut} and still reports $r = \langle C\rangle/\texttt{max\_cut}$.\\[1mm]
\textbf{Statement I:} In that case the reported \texttt{max\_cut} is guaranteed to be at least the true $C_{\max}$.\\[1mm]
\textbf{Statement II:} In that case the reported ratio is at least the true ratio $\langle C\rangle/C_{\max}$.
\begin{itemize}[leftmargin=2.6em, labelsep=0.4em, itemsep=2pt, topsep=3pt]
\item[\optlabel{A}] Both Statement I and Statement II are true
\item[\optlabel{B}] Both Statement I and Statement II are false
\item[\optlabel{C}] Statement I is true, but Statement II is false
\item[\optlabel{D}] Statement I is false, but Statement II is true
\end{itemize}
\end{qitem}
\qsep

\needspace{7\baselineskip}
\begin{qitem}{75}
In the recorded run, the quantum execution time was $40.01$ s for the $80$ objective evaluations (\texttt{cudaq.observe}
calls) of the $225$-gate circuit on 15 qubits. The average time per evaluation and per gate are about
\par\smallskip\noindent\begin{tabularx}{\linewidth}{@{}*{2}{L}@{}}
\optlabel{A}\,$0.010$ s and $44\ \mu$s & \optlabel{B}\,$0.18$ s and $0.8$ ms \\[3pt]
\optlabel{C}\,$0.50$ s and $2.2$ ms & \optlabel{D}\,$40$ s and $0.18$ s
\end{tabularx}
\end{qitem}
\qsep

\needspace{7\baselineskip}
\begin{qitem}{76}
The notebook estimates run time as (number of evaluations) $\times$ (gates per circuit) $\times$ (time per gate), with the time per
gate scaling as $2^n$ for a state vector. Going from the default ($n = 15$, $|E| = 30$, $p = 2$) to $n = 25$ with $|E| = 90$ edges at
the same $p$ and evaluation budget, the run time grows by a factor of about
\par\smallskip\noindent\begin{tabularx}{\linewidth}{@{}*{4}{L}@{}}
\optlabel{A}\,$2.9\times10^{6}$ & \optlabel{B}\,$2.8\times10^{3}$ & \optlabel{C}\,$2.7$ & \optlabel{D}\,$1.0\times10^{3}$
\end{tabularx}
\end{qitem}
\qsep

\needspace{7\baselineskip}
\begin{qitem}{77}
Which statement about QAOA for Max-Cut reflects the current state of knowledge?
\begin{itemize}[leftmargin=2.6em, labelsep=0.4em, itemsep=2pt, topsep=3pt]
\item[\optlabel{A}] No quantum advantage has been demonstrated: the classical Goemans--Williamson algorithm guarantees $0.878$ in polynomial time, and fixed-depth QAOA is limited by locality
\item[\optlabel{B}] The recorded result $r = 0.82$ on 15 nodes proves a quantum speed-up
\item[\optlabel{C}] QAOA solves Max-Cut exactly in polynomial time for any fixed $p$
\item[\optlabel{D}] QAOA with $p = 1$ beats Goemans--Williamson on every 3-regular graph
\end{itemize}
\end{qitem}
\qsep

\needspace{7\baselineskip}
\begin{qitem}{78}
With \texttt{GRAPH = "custom"} and \texttt{CUSTOM\_EDGES = [(0, 1, 1.0), (1, 2, 2.5), (0, 2, -1.0)]} (3 nodes), the maximum cut and
the random-split baseline printed by the notebook are
\par\smallskip\noindent\begin{tabularx}{\linewidth}{@{}*{4}{L}@{}}
\optlabel{A}\,$3.5$ and $2.25$ & \optlabel{B}\,$3.5$ and $1.25$ & \optlabel{C}\,$2.5$ and $1.25$ & \optlabel{D}\,$4.5$ and $2.25$
\end{tabularx}
\end{qitem}
\qsep

\needspace{20\baselineskip}
\begin{passagebox}{Paragraph for Questions 79 and 80}
\emph{Previously recorded run of \texttt{QAOA.ipynb}} (default settings, CUDA-Q 0.16.0, target \texttt{qpp-cpu}):
15 nodes, 30 unit edges ($W = 30$), exact optimum $C_{\max} = 25$ with the four optimal bitstrings listed in A14,
$p = 2$, COBYLA with 80 evaluations, best expected cut $\langle C\rangle = 20.5336$, $4000$ shots. The four most frequent
bitstrings were all optimal:
$\texttt{110010001001101}$ (68 shots), $\texttt{001101110110010}$ (67), $\texttt{110010001011101}$ (56), $\texttt{001101110100010}$ (49);
in total 240 shots were optimal.
\end{passagebox}
\needspace{7\baselineskip}
\begin{qitem}{79}
Estimated from the recorded counts, the probability of sampling the partition described by $\texttt{110010001001101}$
(either this bitstring or its complement) is
\par\smallskip\noindent\begin{tabularx}{\linewidth}{@{}*{4}{L}@{}}
\optlabel{A}\,$0.0338$ & \optlabel{B}\,$0.0600$ & \optlabel{C}\,$0.0170$ & \optlabel{D}\,$0.0263$
\end{tabularx}
\end{qitem}
\qsep

\needspace{7\baselineskip}
\begin{qitem}{80}
A stricter figure of merit measures the improvement over random guessing:
$\eta = \dfrac{\langle C\rangle - W/2}{C_{\max} - W/2}$. For the recorded run, $\eta$ is
\par\smallskip\noindent\begin{tabularx}{\linewidth}{@{}*{4}{L}@{}}
\optlabel{A}\,$0.600$ & \optlabel{B}\,$0.553$ & \optlabel{C}\,$0.221$ & \optlabel{D}\,$0.821$
\end{tabularx}
\end{qitem}
\qsep

\newpage

\begin{center}{\Large\bfseries\color{brand} RESPONSE SHEET}\\[2pt]
{\small Darken exactly ONE bubble per question. Leave all four bubbles empty to skip a question (0 marks).}\end{center}
\textbf{Name:} \hrulefill\quad \textbf{Roll No.:} \rule{35mm}{0.4pt}\par\medskip
\noindent\begin{minipage}[t]{0.235\linewidth}\centering
\textbf{Section A}\\[2pt]
\begin{tabular}{@{}r@{\hspace{5pt}}l@{}}
\textbf{1} & \bub{A}\,\bub{B}\,\bub{C}\,\bub{D} \\[1.6pt]
\textbf{2} & \bub{A}\,\bub{B}\,\bub{C}\,\bub{D} \\[1.6pt]
\textbf{3} & \bub{A}\,\bub{B}\,\bub{C}\,\bub{D} \\[1.6pt]
\textbf{4} & \bub{A}\,\bub{B}\,\bub{C}\,\bub{D} \\[1.6pt]
\textbf{5} & \bub{A}\,\bub{B}\,\bub{C}\,\bub{D} \\[1.6pt]
\textbf{6} & \bub{A}\,\bub{B}\,\bub{C}\,\bub{D} \\[1.6pt]
\textbf{7} & \bub{A}\,\bub{B}\,\bub{C}\,\bub{D} \\[1.6pt]
\textbf{8} & \bub{A}\,\bub{B}\,\bub{C}\,\bub{D} \\[1.6pt]
\textbf{9} & \bub{A}\,\bub{B}\,\bub{C}\,\bub{D} \\[1.6pt]
\textbf{10} & \bub{A}\,\bub{B}\,\bub{C}\,\bub{D} \\[1.6pt]
\textbf{11} & \bub{A}\,\bub{B}\,\bub{C}\,\bub{D} \\[1.6pt]
\textbf{12} & \bub{A}\,\bub{B}\,\bub{C}\,\bub{D} \\[1.6pt]
\textbf{13} & \bub{A}\,\bub{B}\,\bub{C}\,\bub{D} \\[1.6pt]
\textbf{14} & \bub{A}\,\bub{B}\,\bub{C}\,\bub{D} \\[1.6pt]
\textbf{15} & \bub{A}\,\bub{B}\,\bub{C}\,\bub{D} \\[1.6pt]
\textbf{16} & \bub{A}\,\bub{B}\,\bub{C}\,\bub{D} \\[1.6pt]
\textbf{17} & \bub{A}\,\bub{B}\,\bub{C}\,\bub{D} \\[1.6pt]
\textbf{18} & \bub{A}\,\bub{B}\,\bub{C}\,\bub{D} \\[1.6pt]
\textbf{19} & \bub{A}\,\bub{B}\,\bub{C}\,\bub{D} \\[1.6pt]
\textbf{20} & \bub{A}\,\bub{B}\,\bub{C}\,\bub{D} \\[1.6pt]
\end{tabular}
\end{minipage}\hfill\begin{minipage}[t]{0.235\linewidth}\centering
\textbf{Section B}\\[2pt]
\begin{tabular}{@{}r@{\hspace{5pt}}l@{}}
\textbf{21} & \bub{A}\,\bub{B}\,\bub{C}\,\bub{D} \\[1.6pt]
\textbf{22} & \bub{A}\,\bub{B}\,\bub{C}\,\bub{D} \\[1.6pt]
\textbf{23} & \bub{A}\,\bub{B}\,\bub{C}\,\bub{D} \\[1.6pt]
\textbf{24} & \bub{A}\,\bub{B}\,\bub{C}\,\bub{D} \\[1.6pt]
\textbf{25} & \bub{A}\,\bub{B}\,\bub{C}\,\bub{D} \\[1.6pt]
\textbf{26} & \bub{A}\,\bub{B}\,\bub{C}\,\bub{D} \\[1.6pt]
\textbf{27} & \bub{A}\,\bub{B}\,\bub{C}\,\bub{D} \\[1.6pt]
\textbf{28} & \bub{A}\,\bub{B}\,\bub{C}\,\bub{D} \\[1.6pt]
\textbf{29} & \bub{A}\,\bub{B}\,\bub{C}\,\bub{D} \\[1.6pt]
\textbf{30} & \bub{A}\,\bub{B}\,\bub{C}\,\bub{D} \\[1.6pt]
\textbf{31} & \bub{A}\,\bub{B}\,\bub{C}\,\bub{D} \\[1.6pt]
\textbf{32} & \bub{A}\,\bub{B}\,\bub{C}\,\bub{D} \\[1.6pt]
\textbf{33} & \bub{A}\,\bub{B}\,\bub{C}\,\bub{D} \\[1.6pt]
\textbf{34} & \bub{A}\,\bub{B}\,\bub{C}\,\bub{D} \\[1.6pt]
\textbf{35} & \bub{A}\,\bub{B}\,\bub{C}\,\bub{D} \\[1.6pt]
\textbf{36} & \bub{A}\,\bub{B}\,\bub{C}\,\bub{D} \\[1.6pt]
\textbf{37} & \bub{A}\,\bub{B}\,\bub{C}\,\bub{D} \\[1.6pt]
\textbf{38} & \bub{A}\,\bub{B}\,\bub{C}\,\bub{D} \\[1.6pt]
\textbf{39} & \bub{A}\,\bub{B}\,\bub{C}\,\bub{D} \\[1.6pt]
\textbf{40} & \bub{A}\,\bub{B}\,\bub{C}\,\bub{D} \\[1.6pt]
\end{tabular}
\end{minipage}\hfill\begin{minipage}[t]{0.235\linewidth}\centering
\textbf{Section C}\\[2pt]
\begin{tabular}{@{}r@{\hspace{5pt}}l@{}}
\textbf{41} & \bub{A}\,\bub{B}\,\bub{C}\,\bub{D} \\[1.6pt]
\textbf{42} & \bub{A}\,\bub{B}\,\bub{C}\,\bub{D} \\[1.6pt]
\textbf{43} & \bub{A}\,\bub{B}\,\bub{C}\,\bub{D} \\[1.6pt]
\textbf{44} & \bub{A}\,\bub{B}\,\bub{C}\,\bub{D} \\[1.6pt]
\textbf{45} & \bub{A}\,\bub{B}\,\bub{C}\,\bub{D} \\[1.6pt]
\textbf{46} & \bub{A}\,\bub{B}\,\bub{C}\,\bub{D} \\[1.6pt]
\textbf{47} & \bub{A}\,\bub{B}\,\bub{C}\,\bub{D} \\[1.6pt]
\textbf{48} & \bub{A}\,\bub{B}\,\bub{C}\,\bub{D} \\[1.6pt]
\textbf{49} & \bub{A}\,\bub{B}\,\bub{C}\,\bub{D} \\[1.6pt]
\textbf{50} & \bub{A}\,\bub{B}\,\bub{C}\,\bub{D} \\[1.6pt]
\textbf{51} & \bub{A}\,\bub{B}\,\bub{C}\,\bub{D} \\[1.6pt]
\textbf{52} & \bub{A}\,\bub{B}\,\bub{C}\,\bub{D} \\[1.6pt]
\textbf{53} & \bub{A}\,\bub{B}\,\bub{C}\,\bub{D} \\[1.6pt]
\textbf{54} & \bub{A}\,\bub{B}\,\bub{C}\,\bub{D} \\[1.6pt]
\textbf{55} & \bub{A}\,\bub{B}\,\bub{C}\,\bub{D} \\[1.6pt]
\textbf{56} & \bub{A}\,\bub{B}\,\bub{C}\,\bub{D} \\[1.6pt]
\textbf{57} & \bub{A}\,\bub{B}\,\bub{C}\,\bub{D} \\[1.6pt]
\textbf{58} & \bub{A}\,\bub{B}\,\bub{C}\,\bub{D} \\[1.6pt]
\textbf{59} & \bub{A}\,\bub{B}\,\bub{C}\,\bub{D} \\[1.6pt]
\textbf{60} & \bub{A}\,\bub{B}\,\bub{C}\,\bub{D} \\[1.6pt]
\end{tabular}
\end{minipage}\hfill\begin{minipage}[t]{0.235\linewidth}\centering
\textbf{Section D}\\[2pt]
\begin{tabular}{@{}r@{\hspace{5pt}}l@{}}
\textbf{61} & \bub{A}\,\bub{B}\,\bub{C}\,\bub{D} \\[1.6pt]
\textbf{62} & \bub{A}\,\bub{B}\,\bub{C}\,\bub{D} \\[1.6pt]
\textbf{63} & \bub{A}\,\bub{B}\,\bub{C}\,\bub{D} \\[1.6pt]
\textbf{64} & \bub{A}\,\bub{B}\,\bub{C}\,\bub{D} \\[1.6pt]
\textbf{65} & \bub{A}\,\bub{B}\,\bub{C}\,\bub{D} \\[1.6pt]
\textbf{66} & \bub{A}\,\bub{B}\,\bub{C}\,\bub{D} \\[1.6pt]
\textbf{67} & \bub{A}\,\bub{B}\,\bub{C}\,\bub{D} \\[1.6pt]
\textbf{68} & \bub{A}\,\bub{B}\,\bub{C}\,\bub{D} \\[1.6pt]
\textbf{69} & \bub{A}\,\bub{B}\,\bub{C}\,\bub{D} \\[1.6pt]
\textbf{70} & \bub{A}\,\bub{B}\,\bub{C}\,\bub{D} \\[1.6pt]
\textbf{71} & \bub{A}\,\bub{B}\,\bub{C}\,\bub{D} \\[1.6pt]
\textbf{72} & \bub{A}\,\bub{B}\,\bub{C}\,\bub{D} \\[1.6pt]
\textbf{73} & \bub{A}\,\bub{B}\,\bub{C}\,\bub{D} \\[1.6pt]
\textbf{74} & \bub{A}\,\bub{B}\,\bub{C}\,\bub{D} \\[1.6pt]
\textbf{75} & \bub{A}\,\bub{B}\,\bub{C}\,\bub{D} \\[1.6pt]
\textbf{76} & \bub{A}\,\bub{B}\,\bub{C}\,\bub{D} \\[1.6pt]
\textbf{77} & \bub{A}\,\bub{B}\,\bub{C}\,\bub{D} \\[1.6pt]
\textbf{78} & \bub{A}\,\bub{B}\,\bub{C}\,\bub{D} \\[1.6pt]
\textbf{79} & \bub{A}\,\bub{B}\,\bub{C}\,\bub{D} \\[1.6pt]
\textbf{80} & \bub{A}\,\bub{B}\,\bub{C}\,\bub{D} \\[1.6pt]
\end{tabular}
\end{minipage}
\par\bigskip
\begin{center}\textbf{Score computation} (filled in by the evaluator, using the Answer Book)\\[4pt]
\renewcommand{\arraystretch}{1.5}
\begin{tabular}{|c|>{\raggedright\arraybackslash}p{4.6cm}|c|c|c|c|}\hline
\textbf{Sec.} & \textbf{Topic} & \textbf{Attempted} & \textbf{Correct} $N_{\mathrm c}$ & \textbf{Wrong} $N_{\mathrm w}$ & \textbf{Marks} $3N_{\mathrm c} - N_{\mathrm w}$ \\ \hline
A & Mathematics and Graph Theory for QAOA & & & & \\ \hline
B & Quantum Mechanics & & & & \\ \hline
C & QAOA Basics & & & & \\ \hline
D & Advanced QAOA and Implementation & & & & \\ \hline
\multicolumn{5}{|r|}{\textbf{Total (out of 240)}} & \\ \hline
\end{tabular}\end{center}

\end{document}